\section{Post-training operating system y la pila de activos abiertos}

Al final de la entrevista, el centro de la competencia se desplaza de un preentrenamiento único hacia un post-training continuo. Lample sostiene que las arquitecturas de preentrenamiento seguirán siendo parecidas durante bastante tiempo, mientras que el post-training todavía tiene mucho margen para innovar. Esto no significa que el preentrenamiento carezca de importancia, sino que la diferencia entre productos depende cada vez más de la calidad conjunta del environment, el verifier, la generación de datos, el tool use, el serving y el pipeline de experimentación rápida.

\subsection{Por qué post-training no es el nombre de un algoritmo}

\term{post-training} designa en general la adaptación de la conducta posterior al base model, lo que incluye supervised fine-tuning, preference optimization, RL, distillation, tool-use training y evaluation continua. Si cualquiera de estas etapas carece de versiones, métricas y reversión, al sistema le cuesta experimentar con rapidez. Por eso el post-training puede verse como un operating system: en la capa superior recibe las tareas del producto, en la inferior se conecta con el entrenamiento y el serving, y de forma transversal aporta datos, evaluación y gobernanza.

\lecturefigure{14-post-training-os.png}{La competitividad del post-training proviene de un sistema completo que puede iterarse de forma repetible.}{Subtítulos locales 37:10--39:25; diagrama conceptual redibujado.}

\begin{knowledgebox}{Lectura de la figura: por qué la flexibilidad puede superar a la escala por sí sola}
Un post-training online o de alta frecuencia requiere ciclos de experimentación cortos, entornos estables, evaluación automática y una ownership clara. En una organización muy grande cuyo pipeline está disperso y cuyas aprobaciones e interfaces de datos son lentas, el compute adicional no se traduce automáticamente en velocidad de iteración; un equipo pequeño sin disciplina también termina paralizado por experimentos que no pueden reproducirse. La ventaja proviene de «cómputo suficiente + retroalimentación con mucha información + ejecución con poca fricción».
\end{knowledgebox}

\subsection{Cómo entran los modelos abiertos en la pila comercial}

Cuando otros laboratorios publican modelos open-weight potentes, las empresas de soluciones pueden tratarlos como un activo cuyo compute de preentrenamiento ya fue pagado, y luego hacer adaptación al dominio, integración de herramientas y despliegue privado. El valor no desaparece, sino que se desplaza hacia arriba: de la escasez del base model a los datos, el workflow, la confiabilidad y el control.

\lecturefigure{16-open-asset-stack.png}{Los pesos abiertos son una entrada de la pila del producto, no la entrega completa del valor.}{Subtítulos locales 39:30--42:30; materiales oficiales de licencia abierta y despliegue de Mistral 7B.}

\begin{importantbox}{Lista de verificación de ingeniería para la pila de activos abiertos}
\begin{enumerate}
  \item Verificar la license, la ficha del modelo, los datos y las limitaciones conocidas;
  \item Establecer una baseline en el hardware y la workload de destino;
  \item Construir los datos del dominio, la evaluación y los límites de seguridad;
  \item Incorporar herramientas, recuperación, permisos y estado del negocio en pruebas de extremo a extremo;
  \item Preparar pruebas de compatibilidad y una ruta de reversión para las actualizaciones del base model.
\end{enumerate}
\end{importantbox}

\begin{warningbox}{No confundir «reutilizable» con «sin barrera competitiva»}
Los modelos abiertos reducen el costo de obtener capacidades básicas, pero también elevan lo que los clientes esperan en cuanto a migración y control. La verdadera barrera competitiva puede trasladarse a la retroalimentación de alta calidad, los derechos sobre los datos del dominio, la workflow integration, los activos de evaluación, la distribución y la confiabilidad operativa; todo ello sigue exigiendo una inversión de largo plazo.
\end{warningbox}

\subsection{Resumen de la sección}

El post-training es el sistema operativo que conecta las tareas del producto con la conducta del modelo. Los activos abiertos permiten que los equipos reutilicen compute público, pero la diferenciación debe reconstruirse mediante datos, herramientas, despliegue y ciclos cerrados de retroalimentación.
